\documentclass[11pt]{article}
\usepackage[utf8]{inputenc}
\usepackage{amsmath,amssymb}
\usepackage{graphicx}
\usepackage{booktabs}
\usepackage{authblk}
\usepackage[margin=1in]{geometry}
\usepackage[hidelinks]{hyperref}

\title{The Two Structures of Liquid Water from a Parameter-Free Real-Space Model:\\
Isotropy, Directionality, and the Origin of LDL and HDL}

\author[1]{Claes Johnson}
\affil[1]{KTH Royal Institute of Technology, Stockholm, Sweden \\
\texttt{claesjohnson@gmail.com}}
\date{\today \\[0.4em]
\small\textit{Mathematical formulation by the author; implementation and numerical} \\
\small\textit{experiments developed with AI assistance (Claude, Anthropic).}}

\begin{document}
\maketitle

\begin{abstract}
Liquid water is widely modelled as a fluctuating mixture of two local structures: a
low-density liquid (LDL) that is tetrahedral, open and ice-like in short range, and a
high-density liquid (HDL) that is disordered and collapsed with an interstitial fifth
neighbour. Whether real water genuinely realises two liquids is an experimental question;
the mechanistic question underneath it is what molecular feature drives the tetrahedral
LDL structure. We probe this with RealQM, a parameter-free real-space model in which each
electron is a unit-charge component on a non-overlapping spatial domain. In its reduced
form the water oxygen is a single, essentially isotropic electron cloud with two
directional hydrogen donors. This model hydrogen-bonds; it holds an imposed tetrahedral
ice lattice near-perfectly (tetrahedral order $q\approx1$, coordination $\approx4$); it
forms HDL under compression; but it never spontaneously forms LDL, giving instead a single
broad, disordered, under-coordinated liquid at all sizes and densities tested. Attempts to
add directional lone-pair acceptors by angularly splitting the oxygen cloud dissociate the
hydrogen bond. The results converge on one statement: the same isotropy of the oxygen
acceptor that permits hydrogen bonding is what prevents spontaneous tetrahedral ordering.
Hence spontaneous LDL requires directional lone-pair \emph{acceptors}, while HDL is the
default isotropic-packing structure needing no directionality. The model accesses both
structures for clear physical reasons but not their coexistence.
\end{abstract}

\medskip
\noindent\textbf{Keywords:} liquid water; two-state model; LDL/HDL; tetrahedral order;
hydrogen bonding; lone pairs; real-space quantum mechanics; parameter-free.

\section{Introduction}

A long-standing and recently revived picture of liquid water holds that it is, locally, a
fluctuating mixture of two structures~\cite{twostate}. The \emph{low-density liquid} (LDL)
is tetrahedral, four-coordinated and open---the short-range order of ice carried into a
diffusing liquid. The \emph{high-density liquid} (HDL) is disordered and collapsed, with a
fifth neighbour pushed into the first shell. Both are liquids: the split is \emph{within}
the liquid, not the solid--liquid distinction. In deeply supercooled water the two are
hypothesised to become distinct phases separated by a liquid--liquid critical point.

Most evidence for the two-state model comes from empirical water potentials (ST2,
TIP4P/2005 and relatives) tuned to reproduce water, so a two-state signal there can be a
property built into the potential. This motivates a \emph{parameter-free} test. We do not
attempt to settle whether real water has two liquids---a model this reduced cannot, as we
show, even produce the coexistence. Instead we address the mechanistic question that sits
underneath the debate: \textbf{what molecular feature drives the tetrahedral LDL
structure?} A parameter-free model is well suited to this, because whatever it does or does
not produce follows from physics rather than from fitting.

\section{Model and method}

RealQM represents the $N$-electron state as a sum of one-electron components, each carrying
unit charge on a non-overlapping domain of physical space separated by free (Bernoulli)
boundaries; the ground state minimises the standard Coulomb energy over both the components
and the domain partition~\cite{realqm}. Spatial exclusion replaces wave-function
antisymmetry, and cost is set by a fixed grid rather than by the electron count. The
implementation is a browser WebGPU code using a shared field with a per-cell domain label,
which is what allows hundreds of atoms on a laptop.

\paragraph{Reduced water.} Each water is an oxygen carrying a screened valence kernel
(effective charge $+2$, a short core cutoff $r_c$) with a single two-electron cloud, plus
two hydrogens each a unit proton with one electron. The molecular geometry is held near
$104.5^\circ$ by soft harmonic bond-length and bend restraints. The oxygen cloud is thus
essentially \emph{isotropic}: it presents electron density in all directions and carries no
built-in lone-pair lobes. Crucially, the two O--H bonds are \emph{directional}: they aim at
neighbours and provide directional \emph{donors}.

\paragraph{Order parameters.} For every interior oxygen we compute two standard local
descriptors from the running oxygen positions: the tetrahedral order
\[
q = 1 - \tfrac{3}{8}\sum_{j<k}\left(\cos\theta_{jk} + \tfrac13\right)^2
\]
over its four nearest oxygen neighbours ($q\!\to\!1$ perfect tetrahedron, $q\!\to\!0$
disordered), and $d_5$, the distance to the fifth-nearest oxygen (large for the open
tetrahedral LDL, small when an interstitial HDL neighbour intrudes). We also track the
nearest-neighbour distance $d_1$ and the coordination number within $3.5$\,\AA{} (liquid
water $\approx4$, HDL $\approx5$). Histograms are accumulated over molecules and time; a
bimodal $q$ at fixed temperature would be the two-structure signature. To probe density we
use a soft spherical confining wall (a disclosed modelling choice that sets only the average
density, not the local structure); to probe temperature we use overdamped (Brownian)
dynamics with a Langevin thermostat.

\paragraph{Directional-acceptor extension.} To test whether directional lone pairs can be
added, we extended the solver with an angular ``shell split'': the oxygen cloud is divided
into sectors (two or three), each a component confined to its angular wedge about the shared
nucleus, with the wedge frame tracking the molecule's instantaneous orientation. This is the
natural way to give oxygen explicit directional acceptor lobes within the same formalism.

\section{Results}

\subsection{The isotropic water hydrogen-bonds but under-coordinates}
A relaxed dimer of reduced waters in a linear donor--acceptor geometry forms a stable
hydrogen bond, with O--O settling near $3.0$\,\AA{} and holding (the bond is real but on the
weak, long side). In the liquid, however, the same water is markedly under-structured:
confined near liquid density it is disordered, with coordination $\approx2.6$ and no
tetrahedral population. A spherical oxygen accepts a proton from any direction, so it has no
preference for the tetrahedral arrangement. Doubling the system from 64 to 128 waters (a
larger interior core, roughly double the statistics) leaves the $q$ histogram a single broad
hump with no second peak.

\subsection{Ice holds the tetrahedron}
Placed on a diamond-cubic lattice (ice, O--O $=2.75$\,\AA), the same isotropic water holds
the tetrahedral arrangement near-perfectly: $q\approx1$ and coordination $\approx4$. The
reason is instructive: the two directional O--H \emph{donors} supply half of the tetrahedral
scaffold, which---together with the geometric restraints---is enough to keep an
\emph{imposed} lattice metastable, even though the isotropic acceptor cannot help
\emph{choose} that lattice.

\subsection{Melting disorders without densifying}
Heating the ice at fixed volume (melting occurs near $1300$\,K in the model's internal
units) drops the tetrahedral order from $q\approx0.85$ toward $0.55$, but coordination
stays roughly constant near $3.3$. Melting therefore \emph{disorders} the liquid without
\emph{densifying} it: it does not produce HDL. HDL is a density effect, and a fixed volume
cannot supply the required contraction.

\subsection{Compression yields HDL}
Reaching HDL requires the density axis. Compressing the liquid (stepping the confining wall
inward at fixed temperature) forces a fifth neighbour into the first shell: coordination
rises and $d_5$ collapses, while $q$ stays low. The disordered, dense, isotropically-packed
structure is exactly HDL, and the isotropic model produces it readily---no directionality
needed.

\subsection{Carved lone pairs dissociate the hydrogen bond}
The natural remedy---give oxygen explicit directional acceptor lobes---was tried three ways:
a three-fold split (with the higher oxygen kernel this implies) and a charge-neutral
two-hemisphere split biased toward the acceptor side. In every case the dimer
\emph{dissociated} (O--O grew without bound), while the plain isotropic control stayed bound.
The mechanism is clear: in this formalism the hydrogen bond is the donor proton binding to a
single \emph{coherent} acceptor cloud, and putting a free internal boundary through that
cloud destroys the coherent density the proton was holding. The angular-split machinery
itself is sound (validated on a single molecule); it is the water physics that resists.

\begin{table}[t]
\centering
\begin{tabular}{lll}
\toprule
Experiment & Observation & Reading \\
\midrule
Dimer, isotropic O & O--O holds $\sim\!3.0$\,\AA & the cloud binds \\
Liquid, isotropic O & disordered, coord.\ $\sim\!2.6$, no LDL & isotropic $\Rightarrow$ no tetra.\ preference \\
Carved lone pairs & O--O grows, dissociates & splitting breaks the bond \\
Ice (imposed lattice) & $q\!\approx\!1$, coord.\ $\approx4$ & donors hold an imposed tetrahedron \\
Melting (fixed volume) & $q\!:0.85\!\to\!0.55$, coord.\ $\sim\!3.3$ & disorders, cannot densify: not HDL \\
Compression & coord.\ up, $d_5$ collapses & HDL from density, no directionality \\
\bottomrule
\end{tabular}
\caption{Summary of the numerical experiments. Values are exploratory and qualitative;
the pattern, not the third digit, is the result.}
\end{table}

\section{Discussion}

The experiments cohere into a single statement. Textbook water is tetrahedral because oxygen
has both two directional \emph{donors} (the O--H bonds) and two directional \emph{acceptors}
(the lone pairs). The reduced RealQM water has the donors but an \emph{isotropic} acceptor.
That one difference reproduces the whole pattern: the model holds LDL/tetrahedral order when
placed in it (ice, via the donors); it forms HDL when compressed (dense isotropic packing);
and it cannot spontaneously choose LDL, because without a directional acceptor the open
tetrahedral network is never favoured over dense packing.

The sharp conclusion is therefore about the \emph{origin} of the two structures:
\emph{spontaneous LDL requires directional lone-pair acceptors; donor directionality is not
sufficient}. Equivalently, \textbf{HDL is the ``default'' isotropic-packing structure and LDL
is the special, acceptor-stabilised one.} The isotropy that permits hydrogen bonding at all
is the same isotropy that denies the model a directional acceptor---binding and
non-tetrahedrality are two faces of one property.

This also bears on the ``lone pair'' itself. The crisp picture of two equivalent tetrahedral
lone pairs is a localized-orbital model, not a unique observable: photoelectron spectra place
water's two highest occupied orbitals at different energies (not equivalent twins), and the
electron density shows a smeared, acceptor-rich crescent rather than two sharp lobes. What is
robustly observed is four-coordination and an anisotropic, acceptor-rich oxygen density. A
faithful minimal model therefore needs \emph{directional acceptance}, not literal
``rabbit-ears''---and, as the split experiments show, that directionality cannot be obtained
by fragmenting the acceptor cloud without losing the bond.

\section{Relation to DFT: capability versus ablation}

It is fair to ask whether density-functional molecular dynamics (DFT-MD) could run this
study. The question splits in two, with opposite answers.

\emph{Simulating liquid water of this size}---yes, but expensively, and it is a hard case.
Ab-initio DFT-MD of $32$--$128$ waters is a standard, heavily-studied system. However, each
step is a full self-consistent solve; structural statistics require tens to hundreds of
picoseconds, i.e.\ $10^5$--$10^6$ femtosecond-locked steps, placing such runs on
high-performance clusters over days to weeks rather than on a laptop in minutes. Moreover
DFT's water is itself problematic: common semi-local functionals give over-structured,
sluggish water with the wrong density, and even qualitatively reasonable results require
dispersion corrections and carefully chosen functionals. The LDL/HDL regime specifically
lives in \emph{deeply supercooled} water and demands long timescales, large systems and
extensive sampling---which is precisely why the two-state model is studied predominantly
with empirical potentials rather than DFT-MD.

\emph{The ablation experiment}---no, not cleanly. The result here is not an accurate
simulation of water but a \emph{controlled ablation}: the directionality of the oxygen
acceptor was toggled (isotropic cloud versus explicit split lobes) and its causal effect on
tetrahedral ordering observed. DFT provides the fixed electronic structure of real oxygen;
one cannot ask it, in a clean and controlled way, ``what if oxygen had no lone pairs?'' That
single-feature knob---removing a physical ingredient and watching what it was doing---is a
property of a cheap, tunable, parameter-free model, and it is what permits a causal statement
about the driver of LDL.

The honest framing is therefore complementary rather than competitive. DFT would be more
quantitatively accurate for the structure of real water; it cannot cheaply reach the
supercooled two-state regime; and it structurally cannot perform the ablation. The reduced
RealQM water is not competing on accuracy---it is a \emph{conceptual microscope} that
isolates mechanism at laptop cost.

\section{Limitations}

The model is deliberately reduced and the results are qualitative. The water is
under-coordinated relative to real water, so it sits on the disordered/HDL side of the
diagram; its hydrogen bonds are weak; the confining wall imposes an average density and can
induce a density gradient; the dynamics gives diffusive model-time, not calibrated kinetics.
Most importantly, \emph{we never observe a bimodal $q$ at a single temperature}---the
signature of LDL/HDL coexistence. This is not a numerical shortfall but the central result:
the reduced model cannot produce the coexistence, precisely because it lacks the directional
acceptor that would let tetrahedral order compete with dense packing. Testing the two-state
model itself would require a water that both hydrogen-bonds and carries a directional
acceptor---which, we find, angular splitting cannot deliver in this formalism. That is the
open problem this study isolates.

\section{Conclusion}

A parameter-free real-space model with an isotropic oxygen acceptor and directional hydrogen
donors reproduces a coherent account of the two structures of water. It \emph{holds}
tetrahedral (LDL-like) order when that order is imposed, it \emph{forms} HDL under
compression, and it \emph{never spontaneously forms} LDL. The unifying reason is that the
oxygen acceptor is isotropic: this both enables hydrogen bonding and forbids a spontaneous
tetrahedral preference, and attempts to carve in a directional acceptor break the bond. The
study does not decide whether real water is two liquids; it identifies, without any fitted
parameters, the molecular ingredient responsible for the tetrahedral structure---directional
lone-pair acceptance---and reframes HDL as the default isotropic-packing state. Interactive
versions of all experiments are public.

\begin{thebibliography}{9}
\bibitem{twostate} Two-state / two-structure models of liquid water and the low-density /
high-density liquid picture (LDL/HDL), and the hypothesised liquid--liquid transition in
supercooled water. (General reviews of the two-state model.)
\bibitem{realqm} C.\ Johnson, \emph{RealQM: a 3D multiphase continuum computational model for
quantum mechanics}, and companion works on atoms and molecules; public code and interactive
gallery.
\end{thebibliography}

\end{document}
